% Terjemahan appendix2.tex, baris sumber 661--770.
\begin{Exercises}
	\item Misalkan $E$ grup topologis Hausdorff, $M$ subgrup normal hingga dari $E$, dan $E/M$ merupakan grup profinit terhadap topologi hasil bagi. Buktikan bahwa $E$ juga merupakan grup profinit.
	
	\item Definisikan $\hat{\mathbb{N}}$ sebagai himpunan semua pemetaan $n: \{\text{bilangan prima}\} \to \Z_{\geq 0} \sqcup \{+\infty\}$; unsur-unsurnya dapat dinyatakan secara formal sebagai hasil kali $\prod_{p: \text{bilangan prima}} p^{n_p}$. Melalui faktorisasi prima, $\Z_{\geq 1}$ membenam sebagai himpunan bagian dari $\hat{\mathbb{N}}$. Perkalian, pembagi persekutuan terbesar, kelipatan persekutuan terkecil, dan konsep relatif prima di antara ekspresi-ekspresi tersebut didefinisikan dengan cara yang jelas. Misalkan $H$ subgrup tertutup dari grup profinit $G$. Definisikan
	\[ (G:H) := \text{kelipatan persekutuan terkecil dari semua}\; (G/K : H/(H \cap K)) \;\text{dalam $\hat{\mathbb{N}}$}, \]
	di sini $K$ menjelajahi subgrup normal terbuka dari $G$.
	\begin{enumerate}[(i)]
		\item Buktikan bahwa $(G:H)$ juga merupakan kelipatan persekutuan terkecil dari semua $(G:L)$, dengan $L$ menjelajahi subgrup terbuka yang memuat $H$.
		\begin{hint}
			Setiap $L$ semacam itu mesti memuat subgrup terbuka berbentuk $HK$, dengan $K$ seperti di atas.
		\end{hint}
		\item Buktikan bahwa jika $H_2 \subset H_1 \subset G$, maka $(G:H_2) = (G:H_1) (H_1: H_2)$.
		
		\begin{hint}
			Tetapkan $G_K = G/K$, $H_{i, K} = H_i/H_i \cap K$. Diketahui bahwa $(G_K: H_{2, K}) = (G_K: H_{1, K}) (H_{1, K} : H_{2, K})$; ambil kelipatan persekutuan terkecil pada kedua ruas persamaan.
		\end{hint}
		\item Misalkan $H_1 \supset H_2 \supset \cdots$ suatu barisan menurun subgrup-subgrup tertutup, dan $H := \bigcap_i H_i$. Buktikan bahwa $(G:H)$ merupakan kelipatan persekutuan terkecil dari semua $(G:H_i)$.
		% Petunjuk: Pertimbangkan setiap K yang dipilih
		\item Buktikan bahwa $H$ terbuka jika dan hanya jika $(G:H) \in \Z_{\geq 1}$.
	\end{enumerate}

	\item Misalkan $p$ bilangan prima dan $G$ grup profinit. Jika subgrup tertutup $P$ memenuhi
	\begin{inparaenum}
		\item $P$ merupakan grup pro-$p$,
		\item $(G:P) \in \hat{\mathbb{N}}$ relatif prima terhadap $p$,
	\end{inparaenum}
	maka $P$ disebut suatu \emph{subgrup Sylow pro-$p$} dari $G$.
	\begin{enumerate}[(i)]
		\item Buktikan bahwa $G$ selalu memiliki subgrup Sylow pro-$p$.
		\begin{hint}
			Setiap $G/K$ memiliki suatu subgrup Sylow $p$ berupa $P_K$. Terapkan fakta bahwa $\varprojlim$ terfilter dari himpunan-himpunan hingga tak kosong tidak kosong (latihan berikutnya) untuk memilih $P_K$ secara kompatibel untuk setiap $G/K$, lalu tinjau $\varprojlim_K P_K \hookrightarrow G$.
		\end{hint}
		\item Untuk pelengkapan profinit $\hat{\Z} := \varprojlim_{n \geq 1} \Z/n\Z$ dari $\Z$, periksa bahwa grup aditif gelanggang bilangan bulat $p$-adik $\Z_p$ merupakan subgrup Sylow pro-$p$ dari $\hat{\Z}$.
		\item Buktikan bahwa semua subgrup Sylow pro-$p$ saling konjugat.
		\begin{hint}
			Serupa dengan argumen dalam (i), reduksikan ke kasus grup hingga.
		\end{hint}
		\item Buktikan bahwa setiap subgrup pro-$p$ dari $G$ termuat dalam suatu subgrup Sylow pro-$p$.\footnote{Catatan penerjemah: sumber menulis ``subgrup Sylow $p$'' pada butir ini, sedangkan objek yang didefinisikan dan dibahas ialah subgrup Sylow pro-$p$.}
	\end{enumerate}
	
	\item (N.\ Bourbaki \cite[III.58 Théorème 1]{BouE}) Misalkan $(I, \leq)$ poset terfilter tak kosong, dan diberikan funktor $X: I^{\opp} \to \cate{Set}$ yang dituliskan dalam bentuk $(X_i)_{i \in \Obj(I)}$; notasikan pemetaan yang bersesuaian dengan $i \leq j$ sebagai $f_{ij}: X_j \to X_i$. Dengan metode berikut, buktikan bahwa jika setiap $X_i$ merupakan himpunan hingga tak kosong, maka $\varprojlim_i X_i \neq \emptyset$.
	\begin{enumerate}[(i)]
		\item Tinjau keluarga-keluarga himpunan $(A_i)_{i \in \Obj(I)}$ yang memenuhi
		\[ \emptyset \neq A_i \subset X_i, \quad i \leq j \implies f_{ij}(A_j) \subset A_i . \]
		Keluarga-keluarga himpunan ini membentuk suatu himpunan $\Sigma$. Berikan padanya urutan parsial $\preceq$ sebagai berikut: $(A_i)_i \preceq (A'_i)_i$ berarti $A'_i \subset A_i$ berlaku untuk setiap $i$. Gunakan lemma Zorn untuk menunjukkan bahwa $(\Sigma, \preceq)$ memiliki unsur maksimal.
		
		\begin{hint}
			Mudah dilihat bahwa $\Sigma$ tidak kosong. Untuk suatu rantai dalam $(\Sigma, \preceq)$, perlu diperiksa bahwa mengambil irisan pada setiap indeks $i$ tetap memberikan suatu unsur dari $\Sigma$, dan dengan demikian memberikan batas atas bagi rantai tersebut. Sifat ketakkosongan memerlukan syarat bahwa himpunan-himpunannya hingga.
		\end{hint}
		
		\item Buktikan bahwa jika $(A_i)_i$ merupakan unsur maksimal dari $(\Sigma, \preceq)$, maka $f_{ij}(A_j) = A_i$ berlaku untuk setiap $i \leq j$.
		
		\begin{hint}
			Tetapkan $A'_i := \bigcap_{i \leq j} f_{ij}(A_j) \subset A_i$. Masalahnya tereduksi menjadi membuktikan bahwa $(A'_i)_i \in \Sigma$. Sifat ketakkosongan berasal dari syarat bahwa himpunan-himpunannya hingga; selain itu, satu-satunya hal yang masih perlu ditunjukkan adalah $f_{ij}(A'_j) \subset A'_i$. Pertama-tama, perhatikan bahwa $f_{ij}(A'_j) \subset \bigcap_{j \leq k} f_{ik}(A_k)$. Selanjutnya, gunakan sifat terfilter untuk menunjukkan bahwa $\bigcap_{j \leq k} f_{ik}(A_k) = \bigcap_{i \leq h} f_{ih}(A_h) = A'_i$.
		\end{hint}
		
		\item Dengan meneruskan hal di atas, buktikan bahwa setiap $A_i$ dalam unsur maksimal $(A_i)_i$ merupakan himpunan satu titik, dan dengan demikian buktikan bahwa $\varprojlim_i X_i \neq \emptyset$.
		
		\begin{hint}
			Pilih $i \in \Obj(I)$ dan $x_i \in A_i$. Untuk setiap $j$, definisikan
			\[ B_j := \begin{cases}
				A_j \cap f_{ij}^{-1}(x_i), & i \leq j \\
				A_j, & \text{kasus lainnya}.
			\end{cases} \]
			Buktikan bahwa $(B_j)_j \in \Sigma$, sehingga $A_j = B_j$ selalu berlaku, termasuk dalam kasus $i=j$.
		\end{hint}
	\end{enumerate}
	Perhatikan bahwa kasus $(I, \leq) = (\Z_{\geq 0}, \leq)$ dapat ditangani dengan Lema \ref{prop:ML-nonempty}.
	
	\item Jadikan semua perluasan aljabar (atau perluasan hingga) dari medan $F$ sebagai kategori $\cate{ext}_{\mathrm{alg}}(F)$ (atau $\cate{ext}_{\mathrm{f}}(F)$). Tunjukkan bahwa $\cate{ext}_{\mathrm{alg}}(F)$ ekuivalen dengan $\IndC\left(\cate{ext}_{\mathrm{f}}(F)\right)$.
	
	\item Misalkan $A$ suatu koaljabar atas medan $\Bbbk$, yakni koaljabar dalam kategori monoidal $\left( \cate{Vect}(\Bbbk), \otimes_{\Bbbk}\right)$ (Definisi \ref{def:cogebra-monoidal}). Tuliskan $\cated{Comod}A$ untuk kategori komodul kanan $A$ (Definisi \ref{def:comodule-cogebra}), dan tuliskan $\catesubd{Comod}{\mathrm{f}}A$ untuk subkategori penuh yang terdiri atas komodul kanan $A$ berdimensi hingga. Buktikan bahwa $\cated{Comod}A$ ekuivalen dengan $\IndC \left(\catesubd{Comod}{\mathrm{f}}A\right)$.
	
	\begin{hint}
		Serupa dengan kasus ruang vektor dalam Contoh \ref{eg:Ind-Vect}, terapkan Lema \ref{prop:comodule-fd-union}.
	\end{hint}
	
		\item Tinjau subkategori penuh $\mathcal{C}' \subset \mathcal{C}$ dan funktor $F': \mathcal{C}' \to \mathcal{D}$. Asumsikan bahwa
	\begin{compactitem}
		\item $\mathcal{C}'$ kecil, dan semua objeknya kompak dalam $\mathcal{C}$;
		\item setiap objek dari $\mathcal{C}$ dapat dituliskan sebagai $\varinjlim$ kecil terfilter dari objek-objek $\mathcal{C}'$;
		\item $\mathcal{D}$ memiliki $\varinjlim$ kecil terfilter.
	\end{compactitem}
	Buktikan bahwa dalam keadaan ini $F'$ dapat diperluas menjadi $F: \mathcal{C} \to \mathcal{D}$ sedemikian sehingga $F$ mempertahankan $\varinjlim$ kecil terfilter; hingga isomorfisme, $F$ unik.
	% \cite[Tag 0FWY]{stacks}
	
	\begin{hint}
		Pertama-tama, perluas $F'$ menjadi $\widehat{F'}: \IndC(\mathcal{C}') \to \mathcal{D}$ (Definisi--Proposisi \ref{def:Ind-hat-functor}), lalu peroleh $F: \mathcal{C} \to \mathcal{D}$ melalui ekuivalensi dalam Proposisi \ref{prop:recognition-Ind}. Kita juga memerlukan Proposisi \ref{prop:Ind-hat-small}.
		
		Untuk keunikan, perhatikan bahwa jika objek $X$ dari $\mathcal{C}$ dinyatakan sebagai $\varinjlim_i X_i$, dengan $I \to \mathcal{C}'$ suatu funktor dan $I$ suatu kategori kecil terfilter, maka harus berlaku $F(X) \simeq \varinjlim_i F'(X_i)$.
	\end{hint}
	
	\item Misalkan kategori $\mathcal{C}$ memiliki $\varinjlim$ hingga. Buktikan bahwa $\IndC\mathcal{C}$ merupakan kategori $\aleph_0$-presentabel dalam pengertian Definisi \ref{def:presentable-cat}.
	
	\index{ruang Stone@Ruang Stone (Stone space)}
	\index{himpunan profinit@himpunan profinit (profinite set)}
	\item Ruang topologis Hausdorff kompak yang terputus total disebut \emph{ruang Stone}; sesuai konvensi, ruang topologis secara bawaan direalisasikan pada himpunan kecil. Tuliskan $\cate{Stone}$ untuk kategori semua ruang Stone, dan $\cate{FinSet}$ untuk kategori semua himpunan kecil hingga.
	\begin{enumerate}[(i)]
		\item Buktikan bahwa himpunan Cantor dan $\Z_p$ merupakan ruang Stone ($p$ sebarang bilangan prima).
		\item Buktikan bahwa $\cate{Stone}$ ekuivalen dengan $\ProC(\cate{FinSet})$; oleh karena itu, ruang Stone dapat dipandang sebagai himpunan profinit.
		\item Buktikan bahwa grup profinit tidak lain adalah objek grup dalam $\cate{Stone}$.
		\item Buktikan bahwa kategori ruang topologis Hausdorff lokal kompak yang terputus total ekuivalen dengan $\IndC\ProC(\cate{FinSet})$; buktikan bahwa $\Q_p$ merupakan ruang semacam itu.
	\end{enumerate}
	
	\item Periksa rincian Catatan \ref{rem:Ind-k-linear}.
	
	\item Tunjukkan bahwa $\IndC\mathcal{C}$ dapat dicirikan oleh sifat berikut: untuk setiap kategori $\mathcal{D}$ yang memiliki $\varinjlim$ kecil terfilter, funktor yang diinduksi oleh $\mathcal{C} \to \IndC\mathcal{C}$
	\[ \mathrm{Fct}^0(\IndC\mathcal{C}, \mathcal{D}) \to \mathrm{Fct}(\mathcal{C}, \mathcal{D}) \]
	merupakan suatu ekuivalensi. Di sini $\mathrm{Fct}(\cdots)$ menyatakan kategori funktor, sedangkan superskrip $0$ menyatakan subkategori penuh yang terdiri atas funktor-funktor yang mempertahankan $\varinjlim$ kecil terfilter.
\end{Exercises}
